\begin{questions}

\question 计算物体的转动惯量时，可否将物体的质量看作集中在质心处？

\question 设有两个圆盘用密度不同的金属制成，但重量和厚度都相
同，哪个圆盘有较大的转动惯量？

\question 图\ref{fig:10.24} 中是五个物体的横截面，这些截面具有相等的高
度和相等的最大宽度，这些物体又都具有相等的质量。对于与截
面垂直且经过几何中心的转轴来说，哪个物体的转动惯置最大？
哪个最小？

\vspace{1.5em}
\begin{figure}[h]
    \centering
    \includegraphics{figure/fig10.24}
    \caption{}
    \label{fig:10.24}
\end{figure}

\question 本章\ref{sec:10.02} 节中指出，为了描写刚体的一般运动，可将运动
分解为随某选定的基点$ O $的平动加绕$ O $的转动，平动的距离及
\clearpage\noindent
速
% 311.jpg
度随基点的不同选择而有所不同，但转过的角度及角速度是与基
点的选择无关的，这称之为角速度的绝对性。试问：这基点可以
\CJKunderdot{任意}选择吗？可以选在刚体之外的空间某点上吗？

\question 有一均匀的实心圆柱体，沿着一光滑斜面自由落下，问它
滑下时与液下时其末速度是否相等？

\question 在光滑水平面上有一均匀细棒。证明：在力偶作用下，不
管这力偶作用于棒的哪一部位，棒的质心加速度均为零。

\end{questions}
